% !TEX root = main.tex
\subsection{시스템 동작 및 성능 평가}
\subsubsection*{기본 성능을 측정한 전체 공급망 시나리오}
표 3의 비용은 Entry에서 Claim 발급·Exit까지 연결한 실험 1의 대표 Event에서 측정했다.
각 회차는 다음 14개 Event를 실행했다.
\begin{enumerate}
    \item 같은 소유자의 Note 두 개를 Entry한다. 각 State는 $(160,15,115)$다.
    \item 첫 Note를 전량 Transfer하며 운송 탄소 1을 더한다. Recall한 뒤 다시 Transfer하고 Receiver가 Proceed한다.
    \item 두 번째 Note도 Transfer·Proceed한다. Receiver가 두 Note를 Merge해 $(320,30,231)$을 만든다.
    \item Split으로 두 입력을 준비하고 Process $2\to1$로 $(300,30,261)$을 만든다.
    \item 다시 Split한 두 제품 중 하나는 Issue, 다른 하나는 Exit한다.
\end{enumerate}
수치는 $10^9$ scale 정수다. 공식 5회와 warm-up을 포함한 84개 Event의 Raw에 결과 대조 성공이 기록되어 있다.
선택한 대표 위치는 Entry 1, Transfer 3, Proceed 6, Recall 4, Merge 9, Split 12, Process 11, Exit 14, Issue 13이다.

\subsubsection*{추적·동결을 확인한 별도 시나리오}
Entry $A\to(B,C)$, $B\to(D,E)$의 분기 이력은 기본 비용 시나리오와 별개다.
M1-HF에서는 실제 Proof로 해당 이력을 만들고 RPC로 $C,D,E$를 찾아 실제 동결 3건을 제출했다.
이 실행의 기준 블록은 64, 공통 확인 블록은 67이며 결과는 COMPLETE\_AT\_CHECKPOINT다.
\begin{center}\small\setlength{\tabcolsep}{6pt}
    \textbf{표 2. 기능 검증의 근거를 실행별로 구분}\par\smallskip
    \begin{tabular}{@{}p{2.7cm}p{7cm}p{3.8cm}@{}}\toprule
        \textbf{검증 대상} & \textbf{확인한 결과}              & \textbf{근거}        \\ \midrule
        기본 공급망         & 14 Event 연결, 모든 결과 대조 성공     & 실험 1 [D1]          \\
        작은 분기 이력       & C·D·E, 동결 3건, 공통 시점 완료       & M1-HF [D2]         \\
        Claim·Exit 종단  & Claim 2개·Exit 1개 종단 처리       & M1-HF [D2]         \\
        Backward       & Claim에서 원자재 Entry 4개 복원      & M1-HF [D2]         \\
        경합·오류          & 소비 경합은 미완료, 잘못된 chain은 추적 실패 & M1-HF·실험 4 [D2,D3] \\
        \bottomrule\end{tabular}\end{center}

이 결과는 과거 M1의 시험용 verifier 실행이 아니다. M1-HF Raw는 MockVerifiers=false이고, 후속 실험에도 실제 verifier·Proof·receipt가 보존되어 있다.
다만 M1-HF의 20 Event·Process $3\to2$와 기본 실험의 14 Event·Process $2\to1$은 서로 다른 시나리오다.
현재 논문 초안의 두 키 wrapping 방식은 이 한 단계 Master Key 구현과 별도로 대응시켜야 한다.

